% Section 15.7, from lectures/L15/README.md.

\section{Summary}
\label{c15:sec:review}

\needspace{6\baselineskip}
\subsection*{What you should be able to do now}
\begin{itemize}
  \item Implement a shift register deserialising from the bus and discarding stuff bits.
  \item Detect a stuffing error and say exactly what condition constitutes one.
  \item Explain why the receiver samples at \code{sample} and the transmitter shifts at
    \code{bit_done}.
  \item Feed the CRC engine from the receive side without counting stuff bits.
  \item Read somebody else's VHDL and write a review that is useful rather than polite.
\end{itemize}

\needspace{6\baselineskip}
\subsection*{Questions to test yourself with}
\begin{enumerate}
  \item Why does the receiver sample at \code{sample} rather than at \code{bit_done}?
  \item What condition exactly constitutes a stuffing error, and why is it six identical bits that is
    the limit on the receive side when the transmitter counts five?
  \item What happens to the stuff bit when it is correct?
  \item Why may \code{crc15} only see \code{real_bit}, and what would a wrongly fed stuff bit do?
  \item Why is counting received bits not enough to know when a byte is complete?
\end{enumerate}

\needspace{6\baselineskip}
\subsection*{The code review seminar}
About an hour, with the whole class. Every group gets ten minutes and shows a merged pull request
(what it changed, how it was split up, and what the commit messages say), a review the group wrote
(what the reviewer looked for, and what actually changed as a result of the comments), and one thing
that has not worked in the way of working, and what you have done about it.

Pick out, within the group, one merged pull request and one review somebody in the group wrote to
show. It does not have to be the prettiest one; a PR that got justified criticism is worth more
showing than one nobody read.

The seminar is compulsory and is one of the project's two fixed checkpoints. It is not graded
separately, but it feeds into the overall assessment of the way of working described in
appendix~\ref{app:project}.

\needspace{6\baselineskip}
\subsection*{Target}
With \code{rx_shift_reg} finished, all four leaf modules exist, and \code{can_controller} is the only
one left on the CAN side. Groups that are behind should spend the time here rather than starting on
the state machine: it needs all four.

\needspace{6\baselineskip}
\subsection*{Looking ahead}
\Chapref{c16:ch} takes the frame as a state machine, first explained and then built on the transmit
side.
